% =====================================================================
%  Rapport technique - Deep Learning Avance & MLOps - TEK-UP University
%  Sujet : DeepGuard - Detection de visages generes par IA (Deepfake)
%  Auteurs : Sourour Ben Salha & Ghaith Aouissaoui
%  Compilation : pdflatex rapport_technique.tex (2 fois) -- ou Overleaf
% =====================================================================
\documentclass[11pt,a4paper]{article}

\usepackage[utf8]{inputenc}
\usepackage[T1]{fontenc}
\usepackage{lmodern}
\usepackage[french]{babel}
\usepackage[a4paper,top=2.5cm,bottom=2.5cm,left=2.5cm,right=2.5cm,headheight=14pt]{geometry}
\usepackage{microtype}
\usepackage{graphicx}
\usepackage{xcolor}
\usepackage{booktabs,tabularx,multirow}
\usepackage{caption}
\usepackage{enumitem}
\usepackage{fancyhdr}
\usepackage{titlesec}
\usepackage{listings}
\usepackage[most]{tcolorbox}
\usepackage{tikz}
\usetikzlibrary{arrows.meta,positioning,fit,backgrounds,calc,babel}
\usepackage[hidelinks]{hyperref}
\usepackage{float}
\usepackage{needspace}
\usepackage{xstring}

% ---------------- Couleurs (inspirees du logo TEK-UP) ----------------
\definecolor{primaire}{HTML}{1E2A78}
\definecolor{secondaire}{HTML}{2E6DB4}
\definecolor{accent}{HTML}{C0392B}
\definecolor{vert}{HTML}{1E8449}
\definecolor{orange}{HTML}{D35400}
\definecolor{violet}{HTML}{6C3483}
\definecolor{grisclair}{HTML}{F4F6F8}
\definecolor{grismoyen}{HTML}{BFC5CC}
\definecolor{codebg}{HTML}{F7F7F9}

% ---------------- Informations ----------------
\newcommand{\nomprojet}{DeepGuard}
\newcommand{\etudiantA}{Sourour \textsc{Ben Salha}}
\newcommand{\etudiantB}{Ghaith \textsc{Aouissaoui}}
\newcommand{\enseignant}{M. Haythem \textsc{Ghazouani}}

% ---------------- Mise en forme ----------------
\titleformat{\section}{\Large\sffamily\bfseries\color{primaire}}{\thesection}{0.8em}{}[{\color{primaire}\titlerule[0.6pt]}]
\titleformat{\subsection}{\large\sffamily\bfseries\color{secondaire}}{\thesubsection}{0.8em}{}
\titlespacing*{\section}{0pt}{18pt}{10pt}

\pagestyle{fancy}
\fancyhf{}
\fancyhead[L]{\small\sffamily\color{primaire}\textbf{Deep Learning Avancé \& MLOps} | \nomprojet}
\fancyhead[R]{\small\sffamily\color{gray}TEK-UP University}
\fancyfoot[L]{\small\sffamily\color{gray}\etudiantA{} \& \etudiantB}
\fancyfoot[R]{\small\sffamily\color{primaire}\thepage}
\renewcommand{\headrulewidth}{0.4pt}
\renewcommand{\footrulewidth}{0.4pt}

\setlist{itemsep=2pt,topsep=4pt}
\captionsetup{font=small,labelfont={bf,color=primaire}}
\setlength{\parskip}{5pt}
\setlength{\parindent}{0pt}

\lstset{backgroundcolor=\color{codebg},basicstyle=\ttfamily\footnotesize,
  frame=single,rulecolor=\color{grismoyen},breaklines=true,columns=fullflexible,
  keepspaces=true,showstringspaces=false}

\newtcolorbox{encadre}[2][]{enhanced,colback=grisclair,colframe=#2,
  fonttitle=\sffamily\bfseries,coltitle=white,arc=2mm,boxrule=0.8pt,#1}

\tikzset{
  bloc/.style={draw=#1,fill=#1!10,rounded corners=3pt,thick,align=center,
               font=\small\sffamily,minimum height=1cm,text width=2.6cm},
  bloc/.default=secondaire,
  fleche/.style={-{Stealth[length=2.5mm]},thick,color=#1},
  fleche/.default=primaire,
  etiquette/.style={font=\scriptsize\sffamily,color=#1,fill=white,inner sep=1pt},
  etiquette/.default=gray,
}


% ---------------- Elements a completer (rouge) + figures automatiques ----------------
\newcommand{\aremplir}[1]{\textcolor{accent}{\textbf{[à compléter : #1]}}}
\newcommand{\figureou}[4][0.8]{%
  \IfFileExists{figures/#2}{\begin{figure}[H]\centering\includegraphics[width=#1\linewidth]{figures/#2}\caption{#3}\label{#4}\end{figure}}%
  {\begin{encadre}[title=Figure à insérer]{accent}\aremplir{déposer \texttt{\detokenize{#2}} dans \texttt{figures/} : #3}\end{encadre}}}
\newcommand{\ar}{\textcolor{accent}{\textbf{?}}}
\DeclareRobustCommand{\fichier}[1]{\texttt{\StrSubstitute{\detokenize{#1}}{/}{/\allowbreak}}}

\begin{document}

% =====================================================================
%  PAGE DE GARDE
% =====================================================================
\begin{titlepage}
\thispagestyle{empty}
\begin{tikzpicture}[remember picture,overlay]
  \fill[primaire] (current page.north west) rectangle ([yshift=-0.6cm]current page.north east);
  \fill[secondaire] ([yshift=-0.6cm]current page.north west) rectangle ([yshift=-0.75cm]current page.north east);
  \fill[primaire] (current page.south west) rectangle ([yshift=0.6cm]current page.south east);
\end{tikzpicture}
\centering
\vspace*{0.3cm}

\includegraphics[width=7cm]{logo.png}

\vspace{0.8cm}
{\normalsize\sffamily Cycle Ingénieur en Intelligence Artificielle \& Data Science --- 3\textsuperscript{ème} année}

\vspace{1.8cm}
{\large\sffamily\color{secondaire} RAPPORT TECHNIQUE DE PROJET}\\[0.2cm]
{\Large\sffamily\bfseries Deep Learning Avancé \& MLOps}

\vspace{1.4cm}
% ---------------- Sujet entre deux barres ----------------
{\color{primaire}\rule{\linewidth}{1.6pt}}\\[0.5cm]
{\huge\sffamily\bfseries\color{primaire} \nomprojet}\\[0.35cm]
{\LARGE\sffamily\bfseries Détection de Visages Générés par IA\\[0.15cm] (Deepfake : Real vs Fake Faces)}\\[0.4cm]
{\color{primaire}\rule{\linewidth}{1.6pt}}

\vfill

\begin{minipage}[t]{0.48\textwidth}
  \raggedright
  {\sffamily\bfseries\color{primaire} Réalisé par :}\\[0.2cm]
  {\sffamily\large \etudiantA}\\[0.1cm]
  {\sffamily\large \etudiantB}
\end{minipage}%
\hfill
\begin{minipage}[t]{0.48\textwidth}
  \raggedleft
  {\sffamily\bfseries\color{primaire} Encadré par :}\\[0.2cm]
  {\sffamily\large \enseignant}
\end{minipage}

\vspace{1.6cm}
{\sffamily\large Année universitaire 2026 -- 2027}
\vspace{0.8cm}
\end{titlepage}

\tableofcontents
\newpage

% =====================================================================
\section{Introduction et contexte}
% =====================================================================

Les réseaux antagonistes génératifs (GAN) comme StyleGAN produisent des visages photoréalistes de personnes qui
n'existent pas. Ces images (deepfakes) servent à la fraude à l'identité, aux faux profils et à la désinformation, et
l'œil humain ne les distingue plus des vraies photographies.

\nomprojet{} est une application web qui détecte si un visage est réel ou généré par IA, explique visuellement sa
décision et suit un cycle de vie MLOps complet, des données au déploiement conteneurisé. Un seul projet répond à deux
modules : \textbf{Deep Learning Avancé} (axe 4 : Vision Transformers) et \textbf{MLOps}. Chaque module évalue sa
partie du même dépôt Git :
\url{https://github.com/sourourbensalha209-alt/deepguard}.

\begin{encadre}[title=Comment lire ce rapport]{primaire}
\begin{itemize}
  \item \textbf{Partie A (sections 4 à 8)} : modélisation Deep Learning, évaluée par le module Deep Learning Avancé.
  \item \textbf{Partie B (sections 9 à 15)} : cycle de vie MLOps et solution logicielle, évaluée par le module MLOps. Les sections 16 et 17 sont communes.
  \item Les éléments \aremplir{en rouge} dépendent des résultats d'entraînement et sont à compléter avant l'envoi.
\end{itemize}
\end{encadre}

% =====================================================================
\section{Problématique et objectifs}
% =====================================================================

\begin{encadre}[title=Problématique]{accent}
Comment concevoir un système de Deep Learning qui distingue de façon fiable un visage réel d'un visage généré par IA,
explique visuellement sa décision et soit mis à disposition via une application web rapide, reproductible et testée~?
\end{encadre}

\begin{enumerate}
  \item Comparer trois architectures : une baseline CNN, EfficientNet-V2-S (transfer learning) et ViT-B/16.
  \item Expliquer chaque prédiction par une carte de chaleur (CAM pour les CNN, Attention Rollout pour le ViT).
  \item Optimiser l'inférence (export ONNX, quantification INT8) et mesurer le gain de latence.
  \item Servir le modèle par une API REST FastAPI et une interface React (dépôt d'image, webcam, exemples).
  \item Industrialiser : DVC, MLflow (suivi et registre), Docker multi-stage, GitHub Actions, tests, monitoring.
\end{enumerate}

% =====================================================================
\Needspace{8\baselineskip}
\section{Répartition du travail}
% =====================================================================

\begin{table}[H]
\centering\small
\begin{tabularx}{\textwidth}{@{}l>{\raggedright\arraybackslash}X>{\raggedright\arraybackslash}X@{}}
\toprule
\textbf{Membre} & \textbf{Responsabilités principales} & \textbf{Module concerné} \\
\midrule
Sourour \textsc{Ben Salha} & Données, entraînement et comparaison des modèles, explicabilité, DVC, MLflow, export ONNX & Deep Learning Avancé, MLOps \\
Ghaith \textsc{Aouissaoui} & API FastAPI, interface React, Docker, CI/CD GitHub Actions, tests & Deep Learning Avancé (solution full-stack), MLOps \\
\bottomrule
\end{tabularx}
\caption{Répartition des responsabilités. Les deux membres sont solidaires de l'ensemble des livrables ; l'historique Git (commits et Pull Requests) détaille la contribution de chacun.}
\end{table}
\aremplir{ajuster ce tableau à la répartition réelle des tâches}

\newpage
% =====================================================================
\section*{\Huge\color{primaire}Partie A --- Deep Learning Avancé}
\addcontentsline{toc}{section}{\textbf{Partie A --- Deep Learning Avancé}}
% =====================================================================

\section{Données et prétraitement}

\begin{table}[H]
\centering\small
\begin{tabularx}{\textwidth}{@{}l>{\raggedright\arraybackslash}X>{\raggedright\arraybackslash}X@{}}
\toprule
 & \textbf{140k Real and Fake Faces} & \textbf{Real and Fake Face Detection} \\
\midrule
Source & Kaggle (xhlulu) & Kaggle (Yonsei University) \\
Contenu & 70\,000 visages réels (FFHQ) et 70\,000 visages générés (StyleGAN) & 1\,081 visages réels et 960 visages retouchés \\
Résolution & $256\times256$ & $600\times600$ \\
Découpage & train 100\,000 / validation 20\,000 / test 20\,000 & aucun : test de généralisation uniquement \\
\bottomrule
\end{tabularx}
\caption{Jeux de données. Le second contient des visages retouchés et non générés par GAN : une baisse de performance est attendue.}
\end{table}

\textbf{Absence de fuite de données.} Le découpage train / validation / test est fourni par le jeu de données et n'est
jamais modifié. La normalisation utilise les statistiques fixes d'ImageNet (aucun paramètre appris sur le test).
L'augmentation (retournement horizontal, \texttt{ColorJitter}) n'est appliquée qu'à l'entraînement. Les listes de
fichiers sont construites par \fichier{src/data/prepare.py} avec une graine fixe (\texttt{seed: 42}), de façon
équilibrée entre les classes, ce qui rend les sous-ensembles reproductibles.

\textbf{Prétraitement.} Redimensionnement en $224\times224$, conversion en tenseur, normalisation. Les mêmes opérations
sont reproduites à l'identique dans l'API (\fichier{backend/app/inference.py}), après recadrage du visage par un
détecteur OpenCV.

\section{Modèles comparés}

\subsection{Baseline CNN}
Réseau de 4 blocs (convolution $3\times3$, \emph{BatchNorm}, ReLU, \emph{MaxPool}) suivis d'un pooling global et d'une
couche linéaire, entraîné \emph{from scratch}. Il sert de référence.

\subsection{EfficientNet-V2-S (transfer learning)}
Réseau convolutif pré-entraîné sur ImageNet (implémentation \texttt{efficientnetv2\_rw\_s} de la bibliothèque timm,
variante d'EfficientNet-V2-S), dont la tête est remplacée par une couche à 2 classes et dont tous les poids sont affinés.

\subsection{Vision Transformer ViT-B/16}
\textbf{Patchification.} L'image $224\times224$ est découpée en patchs de $16\times16$ pixels, soit une grille de
$14\times14=196$ patchs. Chaque patch est projeté linéairement dans un espace de dimension 768.

\textbf{Embeddings de position.} Un token de classification \texttt{[CLS]} est ajouté, donnant 197 tokens. Des
embeddings de position \emph{appris}, de forme $(1,197,768)$, sont ajoutés aux tokens : sans eux, l'auto-attention
ignorerait l'ordre spatial des patchs.

\textbf{Auto-attention.} Dans chacune des 12 couches (12 têtes), chaque token calcule
\[
  \mathrm{Attention}(Q,K,V)=\mathrm{softmax}\!\left(\frac{QK^{\top}}{\sqrt{d_k}}\right)V ,
\]
ce qui permet de relier des zones éloignées du visage (contour des cheveux, yeux, arrière-plan), là où un noyau
convolutif ne voit que son voisinage. La décision est lue sur le token \texttt{[CLS]}. La grille, le nombre de patchs
et la forme des embeddings sont journalisés dans MLflow par la fonction \texttt{describe\_vit}.

\subsection{Protocole d'entraînement}
\begin{table}[H]
\centering\small
\begin{tabular}{@{}lccc@{}}
\toprule
 & \textbf{Baseline CNN} & \textbf{EfficientNet-V2-S} & \textbf{ViT-B/16} \\
\midrule
Taux d'apprentissage & $10^{-3}$ & $10^{-4}$ & $3\cdot10^{-5}$ \\
Époques & 10 & 3 & 3 \\
Taille de lot & 128 & 64 & 32 \\
\bottomrule
\end{tabular}
\caption{Hyperparamètres utilisés (fichier \fichier{params.yaml}).}
\end{table}
Optimiseur AdamW (décroissance de poids $0{,}01$), planificateur cosinus, précision mixte (AMP) sur GPU, perte
d'entropie croisée. Le meilleur point de contrôle est retenu sur le \textbf{F1 de validation}, avec la classe « fake »
comme classe positive. Graine fixée (42). Chaque modèle est entraîné une seule fois, sur GPU Kaggle (T4~$\times$2), sur 20\,000 images d'entraînement (sur 100\,000 disponibles), avec 4\,000 images de validation et 4\,000 images de test (sur 20\,000 disponibles pour chacun).

\section{Explicabilité}

\textbf{CNN : CAM.} Pour une tête « pooling global + couche linéaire », Grad-CAM est équivalent à CAM : la carte de la
classe $c$ est $\sum_k w_{c,k}\,F_k$, où $F_k$ sont les cartes de caractéristiques de la dernière couche et $w_{c,k}$ les
poids du classifieur. Aucun gradient n'est nécessaire, ce qui permet d'inclure le calcul dans le graphe ONNX.

\textbf{ViT : Attention Rollout.} Pour chaque couche $\ell$, on moyenne l'attention sur les têtes, on ajoute la
connexion résiduelle et on normalise : $\hat A_\ell=\mathrm{norm}(\bar A_\ell+I)$. Le produit
$R=\hat A_L\cdots\hat A_1$ donne, sur la ligne du token \texttt{[CLS]}, l'influence de chaque patch sur la décision.
La carte $14\times14$ est étirée sur la photo. Elle est calculée \emph{dans} le graphe ONNX
(\texttt{ViTWithRollout}) : l'API n'a besoin que d'ONNX Runtime pour prédire et expliquer.

Les cartes indiquent où le modèle a regardé ; elles ne constituent pas une preuve de falsification.

\section{Optimisation de l'inférence}

\textbf{Export ONNX.} Le modèle exporté a une entrée \texttt{image} et deux sorties : \texttt{logits} (classe 0 =
fake) et \texttt{cam} (carte d'explication). Le script vérifie la \textbf{parité} avec PyTorch (écart $<10^{-3}$) et
mesure la latence médiane (p50) et p95 sur CPU, PyTorch contre ONNX Runtime.

\textbf{Quantification INT8.} Une version à poids quantifiés dynamiquement est produite
(\fichier{src/utils/quantize.py}) avec sa taille, sa latence et l'écart maximal de probabilité par rapport au modèle
FP32. Cet écart est mesuré sur des entrées aléatoires : la version INT8 ne doit être déployée qu'après validation sur
le jeu de test.

\section{Résultats et analyse}

Tous les chiffres proviennent d'une exécution unique de chaque modèle (graine 42), sur les sous-ensembles décrits
en section~5.4. Les latences sont mesurées sur CPU, sur la machine du notebook Kaggle, avec une image à la fois.

\subsection{Comparaison des trois modèles}

\begin{table}[H]
\centering\small
\begin{tabular}{@{}lccccc@{}}
\toprule
\textbf{Modèle} & \textbf{Params (M)} & \textbf{Durée (s)} & \textbf{Acc. val} & \textbf{F1 val} & \textbf{AUC val} \\
\midrule
Baseline CNN & 0,39 & 927 & 0,840 & 0,842 & 0,919 \\
EfficientNet-V2-S & 22,15 & 409 & 0,968 & 0,968 & 0,996 \\
ViT-B/16 & 85,80 & 594 & \textbf{0,998} & \textbf{0,998} & \textbf{1,000} \\
\bottomrule
\end{tabular}
\caption{Entraînement et validation (4\,000 images). Durée : temps d'entraînement complet (10 époques pour la baseline, 3 pour les deux autres).}
\end{table}

\begin{table}[H]
\centering\small
\begin{tabular}{@{}lcccc@{}}
\toprule
\textbf{Modèle} & \textbf{Acc. test} & \textbf{F1 test} & \textbf{AUC test} & \textbf{Erreurs / 4\,000} \\
\midrule
Baseline CNN & 0,836 & 0,837 & 0,921 & 656 \\
EfficientNet-V2-S & 0,973 & 0,973 & 0,996 & 108 \\
ViT-B/16 & \textbf{0,996} & \textbf{0,996} & \textbf{1,000} & \textbf{15} \\
\bottomrule
\end{tabular}
\caption{Évaluation sur le jeu de test (4\,000 images, classe positive : « fake »). Le ViT commet environ 7 fois moins d'erreurs qu'EfficientNet et 44 fois moins que la baseline.}
\end{table}

À titre indicatif (approximation binomiale à 95\,\%), l'incertitude sur l'exactitude est d'environ $\pm0{,}2$ point pour le
ViT et $\pm0{,}5$ point pour EfficientNet : l'écart entre les deux modèles (2,3 points) dépasse largement ces marges. Les
chiffres reposent toutefois sur un seul entraînement par modèle : la variabilité entre graines n'a pas été mesurée.

\figureou[0.5]{confusion_matrix.png}{Matrice de confusion du ViT-B/16 sur le jeu de test : 7 vrais visages classés « fake » et 8 faux classés « real ».}{fig:cm}
\figureou[0.5]{roc_curve.png}{Courbe ROC du ViT-B/16 sur le jeu de test (AUC $=1{,}000$ à l'arrondi).}{fig:roc}

\subsection{Généralisation au second jeu de données}

\begin{table}[H]
\centering\small
\begin{tabular}{@{}lccc@{}}
\toprule
\textbf{Modèle} & \textbf{Exactitude} & \textbf{F1 (classe fake)} & \textbf{AUC} \\
\midrule
Baseline CNN & 0,539 & 0,131 & 0,542 \\
EfficientNet-V2-S & 0,524 & 0,040 & 0,479 \\
ViT-B/16 & 0,529 & 0,006 & 0,487 \\
\midrule
Répondre toujours « réel » & 0,530 & 0 & 0,500 \\
\bottomrule
\end{tabular}
\caption{Test de généralisation sur \emph{Real and Fake Face Detection} (1\,081 visages réels, 960 visages retouchés).}
\end{table}

Sur ce jeu, les trois modèles ne font pas mieux qu'un classifieur qui répondrait toujours « réel » (53,0\,\% de visages
réels) : leur AUC est proche de $0{,}5$ et le F1 du ViT sur la classe « fake » est de 0,006, il ne prédit presque jamais
« fake ». La chute est donc complète, et non une simple baisse de performance.

\subsection{Explicabilité}

\figureou[0.98]{heatmaps.png}{Attention Rollout du ViT-B/16 : 4 visages générés (colonnes de gauche) et 4 réels (colonnes de droite), probabilité « fake » en titre.}{fig:heat}
\figureou[0.98]{heatmaps_efficientnet.png}{CAM d'EfficientNet-V2-S sur les mêmes images.}{fig:heat2}

Sur ces 8 images, l'attention maximale du ViT se concentre souvent sur le haut du front et la racine des cheveux, sur le
bord inférieur de l'image (cou, épaules, vêtements) et sur les coins de l'arrière-plan, rarement sur les yeux ou la
bouche. La CAM d'EfficientNet couvre davantage les régions du visage (joues, front, bouche) pour les visages générés,
et les bords ou l'arrière-plan pour les vrais visages. Les cartes sont grossières : $14\times14$ pour le ViT, $7\times7$
pour EfficientNet. Huit exemples ne permettent aucune conclusion générale, et ces cartes indiquent où le modèle a
regardé, pas pourquoi il décide.

\subsection{Optimisation de l'inférence}

\begin{table}[H]
\centering\small
\begin{tabular}{@{}lccccc@{}}
\toprule
\textbf{Modèle} & \textbf{PyTorch} & \textbf{ONNX} & \textbf{INT8} & \textbf{Taille (Mo)} & \textbf{Écart} \\
\midrule
Baseline CNN & 20,3 & 14,5 & 61,1 & 1,56 / 0,40 & 0,0002 \\
EfficientNet-V2-S & 99,9 & 44,0 & 369,5 & 88,4 / 22,9 & 0,0001 \\
ViT-B/16 & 235,9 & 203,9 & \textbf{151,4} & 343,3 / \textbf{86,9} & 0,0329 \\
\bottomrule
\end{tabular}
\caption{Latences médianes p50 en ms (CPU, une image) pour PyTorch, ONNX Runtime (FP32) et INT8 ; taille ONNX FP32 / INT8 ; écart maximal de probabilité entre INT8 et FP32 sur des entrées aléatoires (source : \fichier{metrics/export.json}). L'écart entre PyTorch et ONNX Runtime est inférieur à $2\cdot10^{-5}$ pour les trois modèles.}
\end{table}

L'export ONNX réduit la latence de 1,4 fois (baseline), 2,3 fois (EfficientNet) et 1,16 fois (ViT). La quantification INT8
dynamique n'est utile que pour le ViT : le modèle est 4 fois plus petit (343 $\to$ 87~Mo) et 26\,\% plus rapide que le
FP32 ; pour les deux CNN, elle est \emph{plus lente} (4 fois pour la baseline, 8 fois pour EfficientNet), car la
quantification dynamique cible surtout les produits matriciels. La version INT8 du ViT n'a pas été évaluée sur le jeu de
test : le modèle déployé est le FP32. Dans l'application (Docker sur un ordinateur portable, CPU), le calcul serveur
mesuré va de 0,2 à 2,3~s selon l'image et la charge de la machine, et il atteint 7,6~s au premier appel (démarrage à
froid) ; il inclut la détection du visage et le calcul de la carte d'attention.

\subsection{Robustesse sur une image générée}

Le ViT déployé a été testé sur une seule image générée par StyleGAN, dégradée de plusieurs façons par le script de
robustesse du dépôt (\fichier{scripts/robustness_test.py}). Ce test est qualitatif : une seule image ne permet pas de
conclure.

\begin{table}[H]
\centering\small
\begin{tabular}{@{}lcc@{}}
\toprule
\textbf{Dégradation} & \textbf{Verdict} & \textbf{P(généré)} \\
\midrule
Aucune & FAKE & 99,9\,\% \\
JPEG, qualité 75 / 50 / 25 / 10 & FAKE & 99,6 / 99,9 / 99,9 / 99,5\,\% \\
Flou gaussien, rayon 1 / 2 / 4 & FAKE & 99,9 / 99,8 / 99,7\,\% \\
Bruit, $\sigma=5$ / 15 / 30 & FAKE & 99,7 / 99,7 / 99,8\,\% \\
Réduction à 128 / 64 / 32 px & FAKE / FAKE / \textbf{REAL} & 99,9 / 97,6 / 12,1\,\% \\
Recadrage central 80 / 60\,\% & \textbf{REAL} / \textbf{REAL} & 0,2 / 0,0\,\% \\
\bottomrule
\end{tabular}
\caption{Réponse du ViT à des dégradations d'une image StyleGAN.}
\end{table}

Le modèle résiste à la compression, au flou et au bruit, mais bascule sur « réel » à 32 pixels et après un simple
recadrage central, alors que le visage reste détecté : il dépend du cadrage et de l'échelle du visage, qui sont fixes
dans les données d'entraînement. Ces sorties sont presque toujours proches de 0 ou de 100\,\% : les probabilités ne sont
pas calibrées.

\subsection{Tests qualitatifs sur des images réelles}

Trois essais ont été faits dans l'application, à titre d'illustration. Un vrai visage de face est jugé « réel » à
100\,\% et un visage StyleGAN « généré » à 99\,\%. Une vraie photo retouchée avec un modèle de génération d'images
(ChatGPT) est jugée « réelle » à 100\,\%, le visage étant pourtant détecté. Un gros plan sans visage entier (œil, nez) est
également jugé « réel » à 100\,\%, mais aucun visage n'est détecté et l'image est déformée : l'interface affiche
désormais « Résultat non fiable » dans ce cas.

\subsection{Analyse critique}

\begin{enumerate}
  \item \textbf{Quel modèle gagne ?} Le ViT-B/16 : 99,62\,\% d'exactitude sur le test (15 erreurs sur 4\,000), contre 97,3\,\%
        pour EfficientNet (108 erreurs) et 83,6\,\% pour la baseline (656 erreurs).
  \item \textbf{Son coût est-il justifié ?} Il a 4 fois plus de paramètres qu'EfficientNet (85,8~M contre 22,2~M) et une
        latence ONNX 4,6 fois plus élevée (204~ms contre 44~ms), pour environ 7 fois moins d'erreurs. Son temps
        d'entraînement est comparable (594~s contre 409~s). Pour un usage où 0,2 à 2~s par image sont acceptables, le gain
        justifie le coût ; pour un usage en temps réel, EfficientNet est un compromis plus économique.
  \item \textbf{Que montrent les cartes ?} Voir la section « Explicabilité » : le ViT regarde surtout le haut du front, les
        bords et l'arrière-plan, plus que les yeux et la bouche. Sur 8 images, cela reste une observation, pas une preuve.
  \item \textbf{Pourquoi la performance s'effondre-t-elle sur le second jeu ?} Les « faux » de ce jeu sont des photos
        réelles retouchées, alors que les modèles ont appris sur des visages entièrement générés par StyleGAN. Notre
        hypothèse, non testée, est qu'ils ont appris les traces propres à StyleGAN (et non une notion générale de
        falsification), et que la résolution ($600\times600$ contre $256\times256$) et la chaîne de compression
        différentes de ce jeu déplacent aussi la distribution des images.
  \item \textbf{Quels cas d'erreur ?} Sur le test, 7 vrais visages sont classés « fake » et 8 faux sont classés « real ».
        Ces 15 images n'ont pas été examinées une à une, ce qui est une piste d'analyse.
\end{enumerate}

\begin{encadre}[title=Relecture à faire]{orange}
\aremplir{relire cette analyse avec tes propres observations ; mettre à jour les chiffres si tu relances un entraînement}
\end{encadre}

\newpage
% =====================================================================
\section*{\Huge\color{primaire}Partie B --- MLOps et solution logicielle}
\addcontentsline{toc}{section}{\textbf{Partie B --- MLOps et solution logicielle}}
% =====================================================================

\section{Architecture globale}

L'application est un système découplé, orchestré par Docker Compose : chaque composant a un rôle unique et
communique par des interfaces standard.

\begin{figure}[H]
\centering
\begin{tikzpicture}[node distance=8mm and 9mm]
  \node[bloc=secondaire] (d) {Données\\Kaggle (DVC)};
  \node[bloc=secondaire,right=of d] (p) {\texttt{prepare}};
  \node[bloc=secondaire,right=of p] (t) {\texttt{train}\\(MLflow)};
  \node[bloc=secondaire,right=of t] (e) {\texttt{evaluate}\\\texttt{explain}};
  \node[bloc=vert,below=of e] (x) {Export ONNX\\+ INT8};
  \node[bloc=vert,left=of x] (r) {Registre MLflow\\\texttt{staging} / \texttt{production}};
  \node[bloc=orange,left=of r] (a) {API FastAPI\\ONNX Runtime};
  \node[bloc=violet,left=of a] (f) {Interface\\React};
  \draw[fleche] (d) -- (p); \draw[fleche] (p) -- (t); \draw[fleche] (t) -- (e);
  \draw[fleche] (e) -- (x); \draw[fleche] (x) -- (r); \draw[fleche] (r) -- (a); \draw[fleche] (a) -- (f);
\end{tikzpicture}
\caption{Du jeu de données à l'interface : pipeline reproductible (DVC), suivi (MLflow) et service (FastAPI, React).}
\end{figure}

\begin{table}[H]
\centering\small
\begin{tabularx}{\textwidth}{@{}l>{\raggedright\arraybackslash}X@{}}
\toprule
\textbf{Composant} & \textbf{Rôle} \\
\midrule
\fichier{frontend/} & Interface React, Vite, Tailwind (port 3000) \\
\fichier{backend/} & API FastAPI, validation Pydantic, ONNX Runtime (port 8000) \\
\fichier{src/} & Code Deep Learning : \texttt{data/}, \texttt{models/}, \texttt{utils/} \\
\fichier{dvc.yaml} & Pipeline reproductible \texttt{prepare}, \texttt{train}, \texttt{evaluate}, \texttt{explain}, \texttt{export}, \texttt{register} \\
\fichier{docker-compose.yml} & Lancement complet ; profils \texttt{mlops} (MLflow) et \texttt{monitoring} (Prometheus, Grafana) \\
\fichier{.github/workflows/} & Intégration continue et build des images \\
\bottomrule
\end{tabularx}
\caption{Composants du dépôt.}
\end{table}

\section{Versionnage des données et pipeline reproductible (DVC)}

Le dépôt Git ne contient ni données ni poids : ils sont versionnés par DVC. Le jeu complet (environ 4~Go) reste sur Kaggle ; on
versionne le sous-ensemble de 28\,000 images réellement utilisé, par \texttt{dvc add data/raw} (fichier \fichier{data/raw.dvc}
commité). La graine étant fixée, l'étape \texttt{prepare} retrouve exactement les mêmes listes. Cette étape dépend de
\fichier{data/raw}, et les artefacts sont suivis par le pipeline \fichier{dvc.yaml},
avec un stockage distant DagsHub déclaré dans \fichier{.dvc/config} (les identifiants restent en local, dans
\fichier{.dvc/config.local}, exclu de Git).

\begin{table}[H]
\centering\small
\begin{tabularx}{\textwidth}{@{}l>{\raggedright\arraybackslash}X>{\raggedright\arraybackslash}X@{}}
\toprule
\textbf{Étape} & \textbf{Entrées principales} & \textbf{Sorties} \\
\midrule
\texttt{prepare} & \fichier{data/raw}, \fichier{src/data/prepare.py}, paramètres \texttt{prepare} & listes \fichier{data/processed/*.csv} \\
\texttt{train} & listes, \fichier{src/models/train.py}, paramètres \texttt{train} & \fichier{checkpoints/best.pt}, \fichier{metrics/train.json} \\
\texttt{evaluate} & point de contrôle, listes & matrice de confusion, courbe ROC, \fichier{metrics/test.json} \\
\texttt{explain} & point de contrôle & \fichier{reports/heatmaps.png} \\
\texttt{export} & point de contrôle & \fichier{models/deepguard.onnx} (+ INT8), \fichier{metrics/export.json} \\
\texttt{register} & modèle ONNX, métriques & version du registre MLflow, \fichier{metrics/registry.json} \\
\bottomrule
\end{tabularx}
\caption{Étapes du pipeline DVC. Les hyperparamètres sont dans \fichier{params.yaml} : DVC ne relance que les étapes dont les entrées ont changé.}
\end{table}

\figureou[0.9]{dvc-dag.png}{Graphe du pipeline (\texttt{dvc dag}).}{fig:dag}

L'entraînement lourd s'exécute sur un GPU Kaggle (notebook \fichier{notebooks/kaggle_train.ipynb}) ; ses sorties
sont ensuite enregistrées dans le pipeline par \texttt{dvc commit -f prepare train evaluate explain export} (l'étape \texttt{register} est enregistrée à part, après l'enregistrement du modèle dans MLflow), puis poussées avec \texttt{dvc push}.

\section{Suivi d'expériences et registre de modèles (MLflow)}

Chaque entraînement est un \emph{run} de l'expérience \texttt{deepguard} : paramètres (modèle, hyperparamètres, tailles
d'échantillon, caractéristiques du ViT), métriques par époque (perte, accuracy, F1, ROC-AUC de validation), durée,
\emph{tag} \texttt{git\_sha} du commit. L'évaluation sur le test et le test de généralisation sont journalisés dans
des runs séparés avec leurs figures.

\textbf{Registre.} \texttt{python -m src.models.register} enregistre le modèle ONNX comme nouvelle version de
\texttt{deepguard}. Elle reçoit toujours l'alias \texttt{staging}, et l'alias \texttt{production} si son F1 de
validation est au moins égal à celui de la version actuellement en production (ou sur demande, avec
\texttt{--promote}). MLflow ayant déprécié les « stages » au profit des alias, \texttt{staging} et \texttt{production}
jouent le rôle des statuts Staging et Production. Le script \fichier{scripts/pull_registry_model.py} récupère la
version \texttt{production} dans \fichier{models/}.

\figureou[0.95]{mlflow-runs.png}{Comparaison des trois runs d'entraînement dans MLflow.}{fig:mlflow}
\figureou[0.95]{mlflow-registry.png}{Registre MLflow : versions du modèle et alias \texttt{staging} / \texttt{production}.}{fig:registry}

\section{Backend : API FastAPI}

\begin{table}[H]
\centering\small
\begin{tabularx}{\textwidth}{@{}ll>{\raggedright\arraybackslash}X@{}}
\toprule
\textbf{Méthode} & \textbf{Route} & \textbf{Description} \\
\midrule
GET & \texttt{/health} & État du serveur, GPU disponible, fournisseurs ONNX Runtime, version du modèle, modèle de démonstration ou non \\
GET & \texttt{/model/metadata} & Nom, architecture, classes, métriques de validation et de test, version Git \\
POST & \texttt{/predict} & Image $\rightarrow$ verdict (\texttt{REAL}/\texttt{FAKE}), probabilités, latence serveur \\
POST & \texttt{/predict/visualize} & Idem + calque de chaleur PNG et région analysée (\texttt{face\_box}) \\
GET & \texttt{/metrics} & Métriques Prometheus \\
\bottomrule
\end{tabularx}
\caption{Routes de l'API (documentation Swagger sur \texttt{/docs}).}
\end{table}

\textbf{Robustesse.} Tous les payloads sont validés par des schémas Pydantic. Les erreurs sont sémantiques :
400 si le fichier n'est pas une image valide (format ou contenu), 413 si l'image dépasse 5~Mo ou une taille en pixels
raisonnable, 503 si aucun modèle n'est chargé. L'orientation EXIF des photos de téléphone est corrigée. L'inférence
s'exécute dans un pool de threads pour ne pas bloquer la boucle asynchrone.

\textbf{Modèle de démonstration.} Tant qu'aucun modèle entraîné n'existe, l'API charge un petit modèle ONNX à poids
aléatoires intégré à l'image, afin que toute la chaîne fonctionne ; \texttt{/health} et l'interface le signalent
explicitement, car ses prédictions n'ont aucune valeur.

\textbf{Prétraitement côté serveur.} Un détecteur de visages OpenCV recadre la photo avant l'inférence ; si aucun
visage n'est détecté, l'image entière est analysée et la réponse l'indique (\texttt{face\_detected = false}).

\section{Frontend : interface React}

L'interface (React, Vite, Tailwind) propose trois modes d'ingestion : glisser-déposer, galerie d'exemples en un clic et
webcam. Le résultat affiche le verdict, deux barres de probabilité animées (réel et généré), la latence côté serveur et
côté client, ainsi qu'un curseur de comparaison entre la photo et la carte d'explication superposée à la région
analysée. Un résultat peu sûr (confiance inférieure à 65\,\%) est signalé comme tel plutôt que présenté comme un verdict.
Les composants suivent l'arborescence demandée : \texttt{Dropzone}, \texttt{Webcam}, \texttt{HeatmapCanvas},
\texttt{ProbabilityChart}, \texttt{Samples}, et le client \fichier{services/api.js}.

\figureou[0.95]{ui-accueil.png}{Interface DeepGuard (état initial, modèle de démonstration).}{fig:ui}
\figureou[0.95]{ui-resultat.png}{Interface après analyse d'une image avec le modèle entraîné (verdict, barres, carte d'explication).}{fig:uires}

\Needspace{8\baselineskip}
\section{Conteneurisation et CI/CD}

\textbf{Docker.} Les deux images utilisent des \emph{builds} multi-étapes : l'API construit son environnement Python dans
une première étape puis ne recopie que l'environnement dans une image minimale exécutée par un utilisateur non
privilégié, avec \texttt{HEALTHCHECK} ; le frontend est compilé par Node puis servi par nginx. Le modèle entraîné
n'est pas copié dans l'image : il est monté depuis \fichier{models/} (versionné par DVC). L'application complète se
lance par \texttt{docker-compose up --build}.

\textbf{GitHub Actions.} Deux workflows s'exécutent à chaque push et Pull Request :
\begin{itemize}
  \item \fichier{ci-tests.yml} : lint Python avec \texttt{ruff} (qui remplace \texttt{flake8} et \texttt{black}) et lint du
        frontend avec \texttt{oxlint} (règles React) ; tests de l'API ; build du frontend ; pipeline ML complet sur des
        images synthétiques (CPU).
  \item \fichier{docker-build.yml} : démarrage de toute la stack avec \texttt{docker compose up --wait}, test de fumée
        sur l'API et l'interface, puis construction des images et publication sur GHCR depuis la branche principale.
\end{itemize}

\figureou[0.95]{github-actions.png}{Exécutions des workflows GitHub Actions (tests et build Docker au vert).}{fig:gha}

\Needspace{8\baselineskip}
\section{Tests et monitoring}

\begin{table}[H]
\centering\small
\begin{tabularx}{\textwidth}{@{}l>{\raggedright\arraybackslash}X@{}}
\toprule
\textbf{Niveau} & \textbf{Ce qui est vérifié} \\
\midrule
API (\fichier{backend/tests/}) & 12 tests Pytest : endpoints, formats refusés, fichier corrompu, taille maximale, calque de chaleur, métriques ; modèle ONNX de démonstration réellement chargé \\
Pipeline ML (\fichier{tests/}) & \texttt{prepare} $\rightarrow$ \texttt{train} $\rightarrow$ \texttt{evaluate} $\rightarrow$ explicabilité $\rightarrow$ export ONNX $\rightarrow$ registre sur données synthétiques ; contrat d'entrée et de sorties du modèle ONNX \\
Explicabilité & CAM cohérente avec le CNN ; Attention Rollout du ViT, et parité entre la carte exportée en ONNX et la carte PyTorch \\
Optimisation & quantification INT8 ; règles de promotion du registre (\texttt{staging}, \texttt{production}) \\
Application complète & démarrage de la stack Docker et test de fumée en CI \\
\bottomrule
\end{tabularx}
\caption{Stratégie de tests.}
\end{table}

\textbf{Monitoring.} L'API expose \texttt{/metrics} au format Prometheus : latence et volume par route, nombre de
prédictions par verdict (\texttt{deepguard\_predictions\_total}) et d'images sans visage détecté
(\texttt{deepguard\_no\_face\_total}), utiles pour repérer une dérive (par exemple une explosion du nombre de « fake »).
Prometheus et Grafana se lancent avec le profil \texttt{monitoring}. Un test de charge Locust
(\fichier{tests/load/locustfile.py}) mesure la latence de \texttt{/predict}.

\begin{table}[H]
\centering\small
\begin{tabular}{@{}lcc@{}}
\toprule
\textbf{Scénario Locust} & \textbf{Requêtes/s} & \textbf{Latence p95 (ms)} \\
\midrule
\aremplir{décrire le scénario (utilisateurs, durée)} & \ar & \ar \\
\bottomrule
\end{tabular}
\caption{Résultat du test de charge.}
\end{table}

\Needspace{8\baselineskip}
\section{Limites et éthique}

\begin{itemize}
  \item Le modèle est entraîné sur des visages FFHQ (réels) et StyleGAN (générés) : les images produites par des
        modèles de diffusion récents ne sont pas couvertes.
  \item Une photo de webcam ou téléchargée sur le web (recompression, filtres, cadrage) diffère des images
        d'entraînement : la précision réelle peut être inférieure à celle du jeu de test.
  \item Le détecteur de visages OpenCV (cascade de Haar) est sensible aux visages de profil ou partiellement masqués.
  \item Les cartes d'attention sont des indices et non des preuves.
  \item Un modèle peut se tromper : un verdict ne doit pas être utilisé seul comme preuve.
  \item Les images envoyées sont analysées en mémoire et ne sont pas conservées ; aucune photo de personne n'est
        versionnée sans autorisation.
\end{itemize}

\section{Conclusion et perspectives}

\nomprojet{} livre une chaîne complète : données versionnées, trois architectures comparées, explicabilité exportée avec
le modèle, API typée et testée, interface interactive, images Docker multi-étapes, intégration continue et suivi des
expériences avec registre. Le ViT-B/16 est le meilleur modèle (99,62\,\% d'exactitude sur le test, 15 erreurs sur 4\,000), mais il ne généralise pas au second jeu de données (53\,\%, AUC 0,49) : il détecte les visages générés par StyleGAN, pas les retouches ni les images d'autres générateurs.

Perspectives : entraîner sur des images issues de modèles de diffusion récents ; mesurer la généralisation sur des visages générés par diffusion ; calibrer les probabilités (\emph{temperature scaling}) ; réduire la dépendance au cadrage (une option d'augmentation par recadrages aléatoires, \texttt{augment: scale}, est implémentée et testée en CI mais n'a pas encore été entraînée \aremplir{résultats du ViT « scale » si exécuté}) ; ajouter MixUp/CutMix et un
rééquilibrage de classes, étendre la détection à la vidéo, détecter la dérive de données avec Evidently et déployer sur
le cloud.

\appendix
\section{Reproduire le projet}

\begin{lstlisting}
# Application complete (modele de demonstration si aucun modele entraine)
docker-compose up --build              # interface :3000, API :8000/docs

# Pipeline ML (mini-echantillon pour tester)
pip install torch torchvision
pip install -r requirements-train.txt -r requirements-dev.txt
dvc repro                              # prepare -> train -> evaluate -> explain -> export

# Entrainement complet : notebooks/kaggle_train.ipynb (GPU Kaggle)
# Tests et lint
cd backend && pytest && cd .. && pytest tests -q && ruff check .
\end{lstlisting}

\section{Arborescence du dépôt}

\begin{lstlisting}
deepguard/
|-- .github/workflows/        # ci-tests.yml, docker-build.yml
|-- .dvc/  dvc.yaml  params.yaml
|-- src/{data,models,utils}/  # Deep Learning
|-- backend/                  # FastAPI, tests, Dockerfile, demo_model/
|-- frontend/                 # React : components/, services/
|-- notebooks/kaggle_train.ipynb
|-- tests/  scripts/  monitoring/  mlflow/  docs/rapport/
`-- docker-compose.yml
\end{lstlisting}

\end{document}
